% ============================================================================
% APPENDIX - Cleaned version
% ============================================================================

\appendix

% ============================================================================
\section{Additional Experimental Details}
\label{app:exp-details}

\subsection{Model and Quantization Configuration}
\label{app:model-config}

We use Qwen2.5-7B-Instruct with the following specifications:
\begin{itemize}
    \item Architecture: 28 layers, 28 attention heads per layer (784 heads total)
    \item Hidden dimension: 4096, head dimension: 128
    \item Quantization: GPTQ 4-bit with group size 128
    \item Calibration: 128 samples from C4 dataset
\end{itemize}

\paragraph{Decoding Parameters.}
All experiments use greedy decoding with fixed parameters to ensure controlled 
comparison:
\begin{itemize}
    \item Temperature: 1.0 (greedy)
    \item Top-p/Top-k: disabled
    \item Max generation length: 512 tokens
    \item Repetition penalty: 1.0 (disabled)
\end{itemize}

\subsection{Sample Selection Procedure}
\label{app:sample-selection}

\paragraph{IFEval Benchmark Execution.}
We evaluate both the full-precision (FP16) and GPTQ-4bit quantized models on 
the complete IFEval benchmark using lm-evaluation-harness \cite{eval-harness}. 
This produces per-sample results containing:
\begin{itemize}
    \item Original prompt with explicit instruction constraints
    \item Model-generated response
    \item Constraint-level compliance scores (strict and loose)
    \item Violated constraint types (e.g., language, format, punctuation)
\end{itemize}

\paragraph{Quantization Method Selection.}
Table~\ref{tab:quant-comparison} compares instruction-following 
performance across quantization levels. The 2-bit model produces 
structurally degenerate outputs unsuitable for analysis. The 8-bit 
model slightly outperforms FP16, likely due to regularization effects, 
leaving minimal failure cases for study. We select 4-bit quantization 
as it exhibits measurable but not catastrophic degradation, providing 
a practical regime where instruction-following failures emerge without 
complete model collapse.

\begin{table}[h]
\centering
\small
\begin{tabular}{@{}lcc@{}}
\toprule
\textbf{Quantization} & \textbf{IFEval Avg-4} & \textbf{$\Delta$ from FP16} \\
\midrule
2-bit GPTQ & 0.1353 & $-0.6398$ \\
4-bit GPTQ & 0.7566 & $-0.0185$ \\
8-bit GPTQ & 0.7821 & $+0.0070$ \\
FP16 (baseline) & 0.7751 & $0.0000$ \\
\bottomrule
\end{tabular}
\caption{Instruction-following performance across quantization 
levels on IFEval benchmark. 4-bit quantization provides a balanced 
regime with measurable degradation suitable for mechanistic analysis.}
\label{tab:quant-comparison}
\end{table}

\paragraph{Failure Case Identification.}
From the complete benchmark results, we identify failure cases where:
\begin{enumerate}
    \item FP16 model satisfies all verifiable constraints (pass)
    \item GPTQ-4bit model violates at least one constraint (fail)
    \item Prompt length is between 30-1200 characters
    \item Instructions contain explicit verifiable constraints (language, format, 
          casing, punctuation, structure)
\end{enumerate}

This filtering yields a candidate set of prompts where quantization demonstrably 
degrades instruction-following capability under identical decoding conditions.

\paragraph{Evaluation Protocol and Validation.}
Given the limited pool of clear failure cases and the subtlety of instruction 
violations, we employ a rigorous manual evaluation protocol. Each candidate 
sample is independently assessed by three large language models (GPT-4, 
Claude-3.5-Sonnet, Gemini-1.5-Pro) to verify constraint violations and ensure 
inter-rater reliability. Samples are only included if all three evaluators agree 
on: (1) FP16 satisfies all constraints, and (2) GPTQ-4bit violates at least one 
constraint. This conservative approach ensures high-confidence failure cases but 
further limits the available sample pool.

\paragraph{Representative Sample Selection.}
From the validated failure cases, we manually select 10 representative 
samples to cover diverse failure modes. To ensure these samples capture 
the broader failure distribution, we analyze all candidate failure cases 
and categorize them by violation type. Our selected samples proportionally 
represent the major failure categories observed in the full candidate set:

\begin{itemize}
    \item \textbf{Language switching} (3 samples, $\sim$30\%): Model 
          responds in unintended language (e.g., German, French) despite 
          explicit English instruction. This is the most common failure 
          mode in the candidate set.
    \item \textbf{Format violations} (4 samples, $\sim$40\%): Missing 
          bullet points, paragraph breaks, or structural requirements. 
          Represents the largest category of quantization-induced errors.
    \item \textbf{Constraint omission} (2 samples, $\sim$20\%): Failure 
          to follow specific lexical constraints (forbidden words, casing).
    \item \textbf{Degenerative repetition} (1 sample, $\sim$10\%): 
          Repeated token sequences preventing coherent output.
\end{itemize}

This distribution mirrors the failure pattern observed across the complete 
candidate set, ensuring our detailed analysis on 10 samples reflects the 
typical quantization-induced degradation modes rather than edge cases. The 
selected samples (IFEval IDs: 12, 13, 57, 74, 79, 97, 102, 206, 1075, 1137) 
enable controlled evaluation where each sample has verified ground truth 
(FP16 passes, GPTQ fails), and $\Delta\mathbf{O}^{(\ell,h)}$ can be 
computed exactly by comparing the two models on the same prompt.

\subsection{Hyperparameter Selection}
\label{app:hyperparameters}

\paragraph{Number of Critical Heads ($k$).}
We set $k=15$ based on the distribution of $|\Delta a^{(\ell,h)}|$ values. 
Approximately 15 heads (< 2\% of 784 total) exhibit deviations exceeding 
0.005, forming a natural threshold separating severely degraded heads from 
the majority showing minimal changes.

\paragraph{Stimulation Strength ($\alpha$).}
We evaluate three settings:
\begin{itemize}
    \item $\alpha = 0.0$: Baseline (no intervention)
    \item $\alpha = 5.0$: Moderate stimulation (selected via grid search over 
          $\{1, 3, 5, 7, 10\}$ on 5 held-out failure cases)
    \item $\alpha = 10.0$: Excessive stimulation (to demonstrate over-intervention 
          effects)
\end{itemize}

% ============================================================================
\section{Implementation Details}
\label{app:implementation}

\subsection{Instruction Token Identification}
\label{app:instruction-tokens}

We identify instruction-bearing tokens using regex-based keyword matching on 
the raw prompt text. Example patterns include: \texttt{"two words"}, 
\texttt{"start with"}, \texttt{"json"}, \texttt{"format"}, 
\texttt{"only output.*"}, \texttt{"lowercase"}, etc.

For each matched character span in the text, we map positions to token indices 
using the tokenizer's offset mapping facility, which returns character-level 
boundaries for each token. This yields a set $\mathcal{I}(x) \subset \{1,\ldots,T\}$ 
representing instruction-bearing token positions.

\subsection{Attention Weight Extraction}
\label{app:attention-extraction}

To extract full attention matrices with \texttt{output\_attentions=True}, we 
disable Flash Attention and SDPA optimizations, forcing the model to use eager 
attention implementation. This ensures complete 
$[\text{batch}, \text{heads}, T, T]$ attention tensors are returned for each 
layer, where $T$ is the sequence length.

\subsection{Computing Instruction Attention Score}
\label{app:compute-attention}

Given attention matrix $\mathbf{A}^{(\ell,h)} \in \mathbb{R}^{T \times T}$ 
and instruction token indices $\mathcal{I}(x)$, we compute:

\begin{equation}
a^{(\ell,h)}(x) = \frac{1}{T \cdot |\mathcal{I}(x)|} 
\sum_{t=1}^{T} \sum_{i \in \mathcal{I}(x)} \mathbf{A}^{(\ell,h)}_{t,i}
\end{equation}

This represents the average attention weight from all query positions to 
instruction tokens. We aggregate across failure cases:

\begin{equation}
\Delta a^{(\ell,h)} = \frac{1}{|F|} \sum_{x \in F} 
\left( a^{(\ell,h)}_{\text{FP}}(x) - a^{(\ell,h)}_{\text{Q}}(x) \right)
\end{equation}

Heads are ranked by $|\Delta a^{(\ell,h)}|$ and the top-$k$ form $H^{\star}$.

\subsection{Output Vector Compensation}
\label{app:output-compensation}

\paragraph{Per-Sample Correction.}
For each diagnostic prompt $x \in F$, we cache the output difference:

\begin{equation}
\Delta\mathbf{O}^{(\ell,h)}(x) = 
\mathbf{O}^{(\ell,h)}_{\text{FP}}(x) - \mathbf{O}^{(\ell,h)}_{\text{Q}}(x)
\in \mathbb{R}^{T_x \times d_h}
\end{equation}

where $T_x$ is the sequence length of prompt $x$ and $d_h=128$ is the head 
dimension.

\paragraph{Inference-Time Application.}
We evaluate on the \textit{same} prompts used for diagnosis, ensuring exact 
sequence length matching and isolating the effect of attention-level intervention:

\begin{equation}
\tilde{\mathbf{O}}^{(\ell,h)}(x) = 
\mathbf{O}^{(\ell,h)}_{\text{Q}}(x) + \alpha \cdot \Delta\mathbf{O}^{(\ell,h)}(x)
\end{equation}

This is implemented via PyTorch forward hooks registered on attention head 
modules. During the forward pass, when the quantized model computes 
$\mathbf{O}^{(\ell,h)}_{\text{Q}}(x)$, the hook automatically applies the 
correction before passing the output to subsequent layers. Model weights remain 
unchanged throughout.

\paragraph{Generalization to New Prompts.}
The current implementation requires test prompts to match diagnostic prompts 
for exact sequence length correspondence. For arbitrary prompts, one could 
compute position-averaged corrections:

\begin{equation}
\bar{\Delta\mathbf{O}}^{(\ell,h)} = 
\frac{1}{|F|} \sum_{x \in F} \frac{1}{T_x} \sum_{t=1}^{T_x} 
\Delta\mathbf{O}^{(\ell,h)}(x)[t, :] \in \mathbb{R}^{d_h}
\end{equation}

which can broadcast to any sequence length. We leave this extension to future work.

% ============================================================================
\section{Additional Results}
\label{app:additional-results}

\subsection{Complete Qualitative Comparison}
\label{app:qualitative-table}

Table~\ref{tab:qualitative-full} presents complete outputs for all 10 evaluation 
samples under three stimulation conditions.

\begin{table*}[t]
\centering
\small
\begin{tabular}{@{}cp{4cm}p{4cm}p{4cm}@{}}
\toprule
\textbf{ID} & \textbf{GPTQ ($\alpha=0$)} & \textbf{DAS ($\alpha=5$)} & \textbf{DAS ($\alpha=10$)} \\
\midrule
12  & Degenerative repetition & Correct format \& content & Severe repetition \\
13  & Missing paragraph breaks & Correct structure & Collapsed output \\
57  & Wrong casing \& format & Correct (lowercase + bullets) & Unstable symbols \\
74  & Violates rhyme \& format & Correct limerick & Partial corruption \\
79  & Constraint omission & Fully recovered & Meta commentary \\
97  & \textbf{German output} & \textbf{Correct English} & Off-task generation \\
102 & Missing bullet structure & Correct bullets & Format instability \\
206 & Correct (passed) & Correct (maintained) & Degraded (failed) \\
1075 & Non-JSON format & Valid JSON output & JSON + extra text \\
1137 & French output & Correct structure & Language mixing \\
\bottomrule
\end{tabular}
\caption{Complete qualitative comparison across all 10 evaluation samples. 
Moderate stimulation ($\alpha=5$) recovers 7/8 baseline failures while 
maintaining 2/2 baseline passes. Excessive stimulation ($\alpha=10$) degrades 
all outputs, demonstrating the non-monotonic response to intervention strength.}
\label{tab:qualitative-full}
\end{table*}

\subsection{Top Degraded Attention Heads}
\label{app:per-head-stats}

Table~\ref{tab:top-degraded-heads} lists the 15 most affected heads identified 
by DAS diagnostic procedure:

\begin{table}[htbp]
\centering
\small
\begin{tabular}{@{}cccc@{}}
\toprule
\textbf{Layer} & \textbf{Head} & \textbf{$\Delta a^{(\ell,h)}$} & \textbf{Type} \\
\midrule
21 & 5  & $-0.0121$ & Decrease \\
21 & 0  & $-0.0063$ & Decrease \\
19 & 1  & $-0.0057$ & Decrease \\
14 & 11 & $+0.0052$ & Increase \\
13 & 4  & $-0.0049$ & Decrease \\
23 & 8  & $+0.0047$ & Increase \\
22 & 15 & $-0.0045$ & Decrease \\
20 & 3  & $+0.0044$ & Increase \\
18 & 7  & $-0.0042$ & Decrease \\
24 & 12 & $+0.0041$ & Increase \\
25 & 2  & $-0.0039$ & Decrease \\
21 & 19 & $+0.0038$ & Increase \\
27 & 6  & $-0.0037$ & Decrease \\
26 & 10 & $+0.0036$ & Increase \\
19 & 14 & $-0.0035$ & Decrease \\
\bottomrule
\end{tabular}
\caption{Top 15 attention heads ranked by $|\Delta a^{(\ell,h)}|$. Negative 
values indicate reduced attention to instruction tokens after quantization; 
positive values indicate increased attention. Degraded heads concentrate in 
layers 19-27, consistent with prior work showing upper layers encode high-level 
semantic constraints.}
\label{tab:top-degraded-heads}
\end{table}

\subsection{Attention Deviation Statistics}
\label{app:attention-distribution}

Across all 784 attention heads, the distribution of instruction attention 
deviations is approximately symmetric:
\begin{itemize}
    \item 348 heads (44.4\%) show positive $\Delta a^{(\ell,h)}$ (increased 
          attention to instruction tokens)
    \item 433 heads (55.2\%) show negative $\Delta a^{(\ell,h)}$ (decreased 
          attention to instruction tokens)
    \item Only 15 heads (1.9\%) exhibit $|\Delta a^{(\ell,h)}| > 0.005$
    \item Mean absolute deviation: 0.00118
    \item Standard deviation: 0.00201
\end{itemize}

This symmetric distribution with a small subset of severe outliers supports 
our hypothesis that quantization disrupts a sparse set of critical heads rather 
than causing uniform degradation across the model. The localized nature of 
degradation enables targeted repair through selective head compensation.

% ============================================================================
\section{Computational Resources}
\label{app:compute}

All experiments were conducted on the University of Notre Dame Center for 
Research Computing (CRC) infrastructure:

\paragraph{Hardware.}
\begin{itemize}
    \item GPU: 1$\times$ NVIDIA A10 (24GB GDDR6)
    \item CPU: 32-core Intel Xeon
    \item RAM: 256GB DDR4
\end{itemize}

\paragraph{Runtime.}
\begin{itemize}
    \item IFEval benchmark evaluation (2 models): $\sim$2 hours
    \item Diagnostic phase (attention extraction, 10 prompts): $\sim$15 minutes
    \item Inference with DAS (10 prompts $\times$ 3 $\alpha$ values): $\sim$5 minutes
    \item Total experimental time: < 3 hours
\end{itemize}

\paragraph{Software Environment.}
\begin{itemize}
    \item PyTorch 2.1.0 with CUDA 11.8
    \item Transformers 4.41.0
    \item Auto-GPTQ 0.7.1
    \item lm-evaluation-harness 0.4.2
    \item Python 3.10
\end{itemize}

% ============================================================================
\section{Limitations and Future Directions}
\label{app:limitations}

\subsection{Current Limitations}

\paragraph{Prompt-Specific Correction.}
Our implementation caches $\Delta\mathbf{O}^{(\ell,h)}(x)$ separately for each 
diagnostic prompt, requiring test prompts to match the diagnostic set for 
sequence length compatibility. This controlled setup enables clean before-after 
comparison but limits direct generalization to arbitrary new prompts without 
additional modifications (e.g., position-averaging).

\paragraph{Manual Hyperparameter Selection.}
The stimulation strength $\alpha=5$ is selected via grid search on held-out 
samples and may not transfer optimally across different models, quantization 
methods, or failure types. Developing adaptive or principled selection strategies 
remains an open problem.

\paragraph{Single Model and Quantization Method.}
Our analysis focuses on Qwen2.5-7B with GPTQ-4bit quantization. Whether similar 
localized disruption patterns occur in other model families (LLaMA, Mistral, 
Gemma) or quantization methods (AWQ, SmoothQuant, mixed-precision) requires 
systematic investigation.

\paragraph{Instruction-Following Only.}
We focus exclusively on instruction-following capability. Whether DAS can 
recover other quantization-degraded capabilities (mathematical reasoning, 
factual consistency, safety alignment) through attention-level intervention 
is unexplored.

\paragraph{Evaluation Scale.}
Our detailed analysis examines 10 representative samples to enable thorough 
qualitative assessment. While this reveals clear recovery patterns, large-scale 
evaluation across the full benchmark would strengthen generalizability claims.

\subsection{Future Directions}

\paragraph{Position-Agnostic Correction.}
Developing position-averaged or token-type-specific corrections would enable 
application to arbitrary prompts while maintaining the training-free property.

\paragraph{Automatic Stimulation Strength Selection.}
Investigating lightweight diagnostic signals (e.g., attention entropy, output 
variance) to automatically determine appropriate $\alpha$ values per head or 
per sample.

\paragraph{Multi-Capability Joint Repair.}
Extending DAS to simultaneously recover multiple capabilities by identifying 
capability-specific attention heads and applying orthogonal corrections.

\paragraph{Integration with Quantization Pipeline.}
Incorporating instruction-following diagnostics into quantization calibration 
(e.g., mixed-precision quantization) to preemptively protect critical heads 
during compression.

\paragraph{Theoretical Understanding.}
Developing formal analysis of why instruction-following exhibits localized 
degradation under quantization compared to other capabilities, potentially 
informing more robust quantization-aware training objectives.